\documentclass{article}

% Packages
% \usepackage[utf8]{inputenc}
% \usepackage[T1]{fontenc}
\usepackage{ebutf8}
\usepackage{amsmath}
\usepackage{amsfonts}
% \usepackage{amssymb}
% \usepackage{graphicx}
\usepackage{hyperref}
% \usepackage{cleveref}
% \usepackage{geometry}
% \usepackage{setspace}

% Define example environment
\newtheorem{example}{Example}
\newtheorem{definition}{Definition}
\newtheorem{remark}{Remark}



% Page setup
% \geometry{margin=1in}
% \onehalfspacing

% Title information
\title{A logical framework for dependent sort specifications and their finite models}
\author{Ambroise LAFONT}
\date{\today}

\begin{document}

\maketitle
We present a logical framework for specifying dependent sort specifications (which we call \emph{schemas}) 
and their finite models (which we call \emph{complexes}, as in \cite{leenasubramaniam}).
This is inspired by Uemura's logical framework for SOGATs~\cite{DBLP:journals/mscs/Uemura23}.
\begin{example}
    \label{excat}
    A sort specification for categories can be given as follows:
    \begin{itemize}
        \item $ob: 𝐒𝐞𝐭$
        \item $hom: ob → ob → 𝐒𝐞𝐭$
        \item $id: ∏ x: ob. hom(x, x) → 𝐒𝐞𝐭$
        \item $comp: ∏ x, y, z: ob. hom(x, y) → hom(y, z) → hom(x, z) → 𝐒𝐞𝐭$
    \end{itemize}
Here is an example of a finite model, or complex:
    \begin{itemize}
        \item $a:ob$
        \item $b:ob$
        \item $f:hom(a, a)$
        \item $g:hom(b, a)$
        \item $id_f:id(f)$
    \end{itemize}    
\end{example}

Clearly those specifications look like 
some contexts of some domain specific type theory.
We sketch how to make it formal here.
\paragraph{Disclaimer} I am not sure how original this is;
it may just be a reformulation of Makkai's FOLDS~\cite{MakkaiFirstOL}.

\section{Logical framework}
\subsection{Informal description}
A schema is intuitively a context of (dependent) sort declarations where
each variable has a type of the shape $∏_{a_1:A_1} … ∏_{a_n:A_n} 𝐒𝐞𝐭$.
We may write this type as $Δ → 𝐒𝐞𝐭$, where $Δ$ is the context 
$(a_1 : A_1, …, a_n : A_n)$. Since all sort declarations end with 
"$→ 𝐒𝐞𝐭$", we may as well omit this part. To sum up, a schema is a context 
$x_1 : Δ₁, … x_p : Δₚ$. Now, each $Δᵢ$ describes a complex (or finite model) over the schema
$(x_1 : Δ₁, …, x_{i-1}:Δ_{i-1})$.
\begin{example}
 Let us reformulate the beginning of Example~\ref{excat} with this notation.
 The models of $ob:()$ are sets.
 We have $\hom : (a:ob, b:ob)$. Note that $(a:ob, b:ob)$ can be interpreted as a finite model 
 of $ob:()$, that is, the finite set with two elements $a$ and $b$. A model of $ob,hom$ is a graph.
 Next, we have $id : (x:ob, f:hom(x,x))$. Again, $(x:ob, f:hom(x,x))$ can be interpreted as the finite graph consisting of one looping arrow.
\end{example}

\subsection{Judgements}
\begin{itemize}
    \item $Γ ⊢$ means that $Γ$ is a schema.
\item $Γ| Δ ⊢$ means that $Δ$ is a complex over the schema $Γ$;
\item $Γ | Δ ⊢ A$ means that $A$ is a sort in context $Γ | Δ ⊢$.
 \item $Γ | Δ ⊢ σ ⇒ Δ'$ means that $σ$ is a substitution from the complex $Δ$ to the complex $Δ'$.
 \item $Γ ⊢ t : Δ$ means that the term $t$ has type $Δ$ in $Γ$. In fact, recalling our shortcut notation, we shall actually think of $t$ as a term of type $Δ → 𝐒𝐞𝐭$.

\item $Γ | Δ ⊢ t : A$ means that $t$ is a term of type $A$.

\end{itemize}
Note that we have two distinct judgements for terms: the first one lives in a schema, while the 
second one lives in a complex depending of some schema.

\subsection{Typing rules}
There are the obvious context extension rules for schemas and complexes.
\[
\dfrac{Γ|Δ⊢}{Γ,x:Δ}
\qquad 
\dfrac{Γ|Δ⊢A}{Γ|Δ,x:A ⊢}
\]
We have the obvious variable rules for schemas and complexes
\[
\dfrac{(x : Δ) ∈ Γ}{Γ⊢x:Δ}
\qquad 
\dfrac{(x:A)∈Δ}{Δ⊢x:A}
\]
In fact this is the only rule for terms! 
In other words, \textbf{all terms are variables}.

We omit the rules for building substitutions between complexes (substitutions are suitable lists of terms) which are standard.
The most important rule is the sort formation rule: 
\[
\dfrac{Γ⊢t : Δ'\quad Γ| Δ⊢σ ⇒ Δ'}{Γ|Δ⊢t(σ)}
\]
This rule makes sense if we think of $t:Δ'$ as $t:Δ' → 𝐒𝐞𝐭$.
\begin{remark}
    The notational convention consisting in 
    omitting the part "$→𝐒𝐞𝐭$" may look confusing.
    The main benefit is that we don't have to give $𝐒𝐞𝐭$ a meaning, 
    thus it makes the semantics of the logical framework simpler.
\end{remark}
\subsection{Rocq formalisation}
Matthieu Piquerez has a formalization of the logical framework
based on raw syntax based on lists of natural numbers, with 
\begin{itemize}
    \item well-formed boolean predicates;
    \item a Ltac DSL as an interface to build contexts;
    \item pushouts of complexes (morphisms are substitutions between complexes).
\end{itemize} 
\subsection{
Semantics
}
It is well-known that schemas correspond to certain kind of finite categories called \emph{simple}.
Then, the category of models of a schema is equivalent (but not isomorphic) to a 
category of functors from the corresponding simple category to $𝐒𝐞𝐭$.
\begin{example}
    A model of $ob:(), hom: (a:ob,b:ob)$ is a graph, that is: a set of vertices $ob$ and a set of edges $hom$ indexed by $ob × ob$.
    The corresponding simple category is $⇉$, and a functor to $𝐒𝐞𝐭$ consists of 
    a set of objects $ob$, a set of arrows $hom$, and two maps $dom, codom : hom → ob$.
\end{example}
\begin{remark}
    As argued before, a complex can be thought of as a finite model of a schema, or a "finite" functor $F$ from a simple category $C$ to $𝐒𝐞𝐭$. The length of the complex is exactly the cardinal 
    of $∑_{c∈C} F(c)$. For example, the complex $a:ob, b:ob, f:hom(a,b)$ is of size 3.
\end{remark}

A schema $Γ$ can be interpreted as:
\begin{itemize}
    \item a simple category $C_Γ$;
    \item a category of models $⟦Γ ⟧$, which is equivalent to $[C,𝐒𝐞𝐭]$,
\end{itemize}
A complex $Γ|Δ$ can be interpreted as an object $⟦Δ⟧$ of the category of models of $Γ$, or more abstractly, as a functor\footnote{This is a route followed by \cite{DBLP:journals/corr/AltenkirchCDF16}} from $Γ$ to $𝐒𝐞𝐭$. Indeed, any object induces such a functor $yΔ$ by the Yoneda embedding.


The context $⟦Γ,x:Δ ⟧$ can be interpreted 
as the category $(γ:⟦Δ⟧) × 𝐒𝐞𝐭^{\hom(⟦Δ⟧,γ)}$.

Again, a type $Γ | Δ ⊢ A$ can be interpreted as an object $⟦A⟧$ of the coslice category $⟦Δ⟧ / ⟦Γ⟧$.
Then, we can define  $⟦Γ|Δ,x:A^+⟧$ as the coslice category $⟦A ⟧ / ⟦Γ⟧$.


The interpretation of terms is then derived: they are sections of the projections from extended contexts to the base context.

Substitution between schemas must be interpreted as adjunctions between categories of models,
so that substituted complexes are interpreted by applying the left adjoint to the original complex.
The most interesting part is the interpretation of the application rule (TODO: provide some details.).

This intepretation yields a model of the logical framework. This model which should actually be equivalent to the syntactic model. In particular, any simple category should 
induce a syntactic schema, and any finite model should induce a syntactic complex.
\section{A single-context logical framework}

In the provided semantics, there is some redundancy, in the sense that the interpretation 
of a complex is very similar to the interpretation of a type. 
In this section, we present a single-context logical framework, with simpler semantics, where 
contexts are no longer of the shape $Γ|Δ$, but just some $Γ$.
Contrary to the previous logical framework which cleanly separates schemas and complexes,
the context can now interleave sort declarations and 
complex declarations, so that the following context is valid:
$ob:𝐒𝐞𝐭, a:ob, hom : ob → ob → 𝐒𝐞𝐭$.
The interest of this single version is that:
\begin{itemize}
    \item The semantics of this logical framework is simpler in some respect;
    \item Any model of this logical framework induces a model 
    of the previous one\footnote{In particular, by initiality/recursion, we have a syntactic translation from the two-context logical framework to the single-context logical framework.}.
\end{itemize}
Therefore, we can exhibit a model of the original 2-context version from a model of the 
1-context version.


\subsection{Polarities}
A context is a list of variable declarations
$x_1 : A_1^{+/-}, …, x_n : A_n^{+/-}$, that is,
each variable $x_i$ has an associated type $A_i$ and a "polarity" $+$ or $-$. 
A negative declaration $x:A^-$ means that we should think of $x$ as of type $A⁺ → 𝐒𝐞𝐭$, so that we 
have the following application rule 
\[
\dfrac{Γ⊢t:A^- \quad Γ⊢δ:A^+}{Γ⊢t(δ)}
\]
We have two extension rules for contexts, depending on the chosen polarity
\[
\dfrac{Γ⊢A}{Γ,x:A^+ ⊢}
\qquad 
\dfrac{Γ⊢A}{Γ,x:A^- ⊢}
\]
Note that we have two term judgements $Γ⊢t:A^+$ and $Γ⊢t:A^-$, depending on the polarity.
\subsection{Type formers}
We assume a unit type and a sigma type with domain in a positive type. 
\[
\dfrac{}{Γ ⊢ 1}
\qquad 
\dfrac{Γ ⊢ A^+ \quad Γ, a:A^+ ⊢ B}{Γ ⊢ ∑_{x:A} B}
\]
Such that we have projections and pairing as usual for $(∑_{x:A})^+$, and a unique term
$Γ ⊢ * : 1⁺$.

\subsection{Example: categories}
The context for categories is: 
$
ob:1^-, hom: (ob × ob)⁻ 
, id : (∑_{x:ob} hom(x,x))^-
$

\subsection{Semantics}
The semantics is similar to the 2-context version, although we can no longer directly interpret 
a context as a simple category since we can interleave positive and negative declarations.

\bibliographystyle{IEEEtran}
\bibliography{bib}

\end{document}